\subsection{Multi-LoRA Scheduler}
\label{sec:scheduler}

\Name not only reduces runtime overhead through fused kernels but also improves end-to-end throughput by scheduling multiple LoRA fine-tuning jobs together.
This is achieved by grouping adapters and adaptively batching their samples to balance GPU load and minimize distributed parallelism overhead.
Figure~\ref{fig:system-design-fig-4-scheduler} shows the overall process. 
Adapters are grouped based on sequence length statistics (top), and adapters in the same group are trained jointly.
For each group, we aggregate samples into global batches and pack each one into microbatches using a two-stage MILP-based optimization (middle).
Once microbatches are generated in parallel with multiprocessing, a final merge pass combines underfilled microbatches across global batches when data dependencies allow (bottom).

\myparagraph{Granularity}
Due to data dependencies between consecutive global batches, our scheduling operates at the granularity of individual global batches. Each adapter's dataset is divided into global batches based on the user-specified global batch size. We then aggregate all samples belonging to the same global batch index across adapters and pack them into multiple microbatches.

\begin{figure}[t]
    \centering
    \includegraphics[width=0.93\linewidth]{figures/system-design-fig-4-scheduler.pdf}
    \caption{Multi-LoRA adapter scheduling workflow. Top: Adapter grouping by sequence length statistics. Middle: Two-stage MILP optimization for microbatch creation. Bottom: Cross-batch merging of underfilled microbatches.}
    \label{fig:system-design-fig-4-scheduler}
\end{figure}

\myparagraph{Bubble Lemma \& Adapter Grouping}
\Name first groups LoRA adapters before batching samples to reduce scheduling complexity.
% While improving load balance through batching and merging, we must preserve pipeline consistency.
In pipeline parallelism with $S$ stages, a sample's backward pass begins only after $S - 1$ other microbatches complete forward passes.
To maintain data dependencies, we define the \textit{bubble lemma}:
for adapter $i$, if sample $s$ from global batch $j$ is committed at microbatch $k$, no sample from batch $j$+1 of the same adapter can be committed before microbatch $k$+$S$-1 (after sample $s$'s backward pass completes).
Without constraints, adaptive batching might scatter samples from consecutive batches across microbatches.
The tail of batch $j$ could conflict with batch $j$+1's head, causing incorrect execution.
We resolve this conflict by grouping adapters with strict ordering between groups, while allowing flexible merging within each group.
This ensures batches from the same adapter are separated by microbatches from other groups, satisfying the bubble condition.
For load balance within groups, we use head-tail pairing, sorting adapters by mean token length, and pairing short-sequence adapters with long ones.
This grouping balances constraints with flexibility and forms our data packing foundation.

\begin{algorithm}[t]
    \footnotesize
    \caption{Data Batching \& Merging (Per Group)}
    \label{alg:multi-lora-data-batching}
    \DontPrintSemicolon
    \KwData{Adapters with grouped samples, token capacity $C$, timeout $t$}
    \KwResult{Scheduled microbatches satisfying pipeline constraints}
    
    \ForEach{global batch $b$ in parallel}{
        \tcp{Greedy fallback as baseline}
        $(B_g, \{m_i^g\}) \gets \texttt{GreedyPacking}(b, C)$
        
        \tcp{Stage 1: minimize number of microbatches}
        $B^* \gets \texttt{MILP\_MinBins}(b, C, \text{timeout}=t)$ \;
        \If{$B^* \geq B_g$}{
            $B^* \gets B_g$
            }
            
            \tcp{Stage 2: minimize smallest bin tokens}
            $\{m_i\}_{i=1}^{B^*} \gets \texttt{MILP\_MinSmallestBin}(b, B^*, C, \text{timeout}=t)$ \;
            \If{$B^* = B_g$ \textbf{and} $\min_i m_i \geq \min_i m_i^g$}{
                \Return $\texttt{GreedyPacking}(b, C)$
                }
                }
                
                % \tcp{Merge underfilled microbatches across batches}
                \ForEach{consecutive batch pairs $(b, b{+}1)$}{
                    Shift tokens from $b{+}1$ into $b$ if bubble lemma is preserved
                    }
                    
                    \texttt{VerifyAndFix(schedule)} \tcp*{Insert no-ops where needed}
                    
    \KwRet Scheduled microbatches
\end{algorithm}

\myparagraph{Data Batching with Two-Stage MILP}
After adapters are grouped, we solve a bin-packing problem to batch samples into microbatches, each constrained by a fixed token capacity.
Our goal is to reduce both the total number of microbatches and the impact of underfilled ones, which affect load balance and pipeline utilization. Specifically, we aim to (i) minimize the number of microbatches needed to pack all samples, and (ii) make the smallest microbatch as empty as possible to enable better merging in later stages.

We address the bin-packing problem using a two-stage mixed-integer linear programming (MILP) formulation (see Algorithm~\ref{alg:multi-lora-data-batching}, lines 3-7).
For notational brevity, let $P$ denote the padding multiple, which is a user-specified parameter to pad the sequence length belonging to the same adapter to a multiple of $P$ (e.g., 64 or 128). Let $x_{s,b} \in \{0,1\}$ denote whether sample $s$ is assigned to bin $b$, $k_{a,b} \in \mathbb{N}$ be the padded multiples contributed by adapter $a$ in bin $b$, and $z_b \in \{0,1\}$ indicate whether bin $b$ is used.  In the first stage, we solve:
\begin{equation}
\begin{aligned}
    \argmin_{x_{s,b}, k_{a,b}, z_b} & \quad \sum_{b=1}^B z_b \\
    \text{s.t.} \quad & \quad z_{b+1} \leq z_b \quad \forall b < B \\
    &\quad \sum_{b=1}^B x_{s,b} = 1 \quad \forall s \in \text{samples} \\
    &\quad \sum_{s \in \text{adapter}(a)} \text{len}(s) \cdot x_{s,b} \leq k_{a,b} \cdot P \quad \forall a,b \\
    &\quad  z_b \leq \sum_{a} k_{a,b} \cdot P \leq \text{capacity} \cdot z_b \quad \forall b
\end{aligned}
\end{equation}
The constraints ensure used bins are contiguous from the start, each sample is assigned to exactly one bin, adapter-specific token counts respect padding multiples, and bin capacity is not exceeded.

With the optimal number of bins $B^{*}$ from the first stage, the second stage fixes $B = B^{*}$ and minimizes the smallest total token count among all bins. The second stage solves:
\begin{equation}
\begin{aligned}
    \argmin_{x_{s,b}, k_{a,b}} & \quad \min_{b \in [1, B^*]} \sum_{a} k_{a,b} \cdot P \\
    \text{s.t.} \quad & \quad \sum_{b=1}^{B^*} x_{s,b} = 1 \quad \forall s \in \text{samples} \\
    &\quad \sum_{s \in \text{adapter}(a)} \text{len}(s) \cdot x_{s,b} \leq k_{a,b} \cdot P \quad \forall a,b \\
    &\quad \sum_{a} k_{a,b} \cdot P \leq \text{capacity} \quad \forall b \\
\end{aligned}
\end{equation}
This optimization problem can be reformulated into an MILP problem. It minimizes the smallest bin size, which leaves more space in the least-full microbatch for potential merging in later stages, thereby reducing pipeline stalls. 

To improve runtime efficiency, we implement two techniques (also reflected in Algorithm~\ref{alg:multi-lora-data-batching}). First, we set a timeout on the MILP solver and fall back to a greedy bin-packing algorithm if the solver takes too long (lines 2, 5, and 9). Second, since global batches are independent, we parallelize the bin-packing optimization across batches using multiprocessing (line 1), which allows us to efficiently schedule all training data while balancing load and reducing microbatch count.

\myparagraph{Merging \& Verification}
After microbatch packing, the final microbatch in a global batch may be underfilled, reducing GPU efficiency and increasing pipeline bubbles. To mitigate this issue, we apply a greedy merge pass that shifts tokens from the next global batch into the current batch's final microbatch (as shown in Figure~\ref{fig:system-design-fig-4-scheduler} bottom), as long as capacity and the bubble lemma are satisfied (Algorithm~\ref{alg:multi-lora-data-batching}, lines 12-14). We then perform a verification step to ensure no constraint is violated. If any bubble condition is unmet, we insert no-op microbatches to restore correctness and preserve pipeline consistency (line 15).

\myparagraph{Parallelism Profiler}
The scheduler requires token capacity as input, which depends on the parallelism strategy.
\Name assumes that effective scheduling keeps tokens within each microbatch close to the token capacity.
Since token capacity and parallelism strategies are orthogonal to scheduling, they should be tuned outside the scheduler using automatic parallelization techniques~\cite{zheng2022alpa,flexflow,liu2024aceso,zhu2025mist}.
We implement a lightweight profiler that directly benchmarks runtime under different model parallelism configurations with fixed-length inputs and collects throughput. 
We choose the best-performing configuration and pass its token capacity to the data batching stage, ensuring that microbatch packing aligns with the system's performance characteristics.

    